\documentclass[10pt,conference]{IEEEtran}
\usepackage[T1]{fontenc}
\usepackage{times}
\usepackage{amsmath,amssymb}
\usepackage{booktabs}
\usepackage{listings}
\usepackage{xcolor}
\usepackage[hidelinks]{hyperref}
\usepackage{url}
\hypersetup{pdftitle={The Multimodal Mind: Perceive, Imagine, Act},pdfauthor={Siddharth Choudhary}}

\lstset{
  basicstyle=\ttfamily\footnotesize,
  breaklines=true,
  columns=fullflexible,
  keepspaces=true,
  xleftmargin=1em,
  frame=leftline,
  framesep=6pt,
  rulecolor=\color{gray!50},
  aboveskip=6pt, belowskip=6pt,
}

\title{The Multimodal Mind: Perceive, Imagine, Act}
\author{\IEEEauthorblockN{Siddharth Choudhary}}

\begin{document}
\maketitle

\begin{abstract}
The next frontier in AI is a single model that understands the physical world
well enough to imagine its futures and act to shape them. The path runs through a
natively multimodal foundation model---one that perceives, reasons, and generates
across all modalities (image, video, text, speech, action)---organized around a
three-stage cognitive loop: Perceive $\rightarrow$ Imagine $\rightarrow$ Act.
Biological intelligence follows that loop without thinking about it, yet current
AI breaks it into disconnected pipelines, each trained by a different team on a
different dataset. This essay argues for building it as one model: a
mixture-of-transformers backbone that runs the loop, a staged pretraining
curriculum that climbs from understanding to generation to action, and a
self-play flywheel that grounds learning in physics. The unifying observation is
that every major AI capability challenge---controllable video, computer use,
robotics, world simulation, autonomous driving---is a special case of multimodal
conditional generation, so one architecture subsumes them all.\footnote{\url{https://itzsid.github.io/publications/multimodal-mind.html}}
\end{abstract}

\section{The Thesis}
The next frontier in AI is a single model that understands the physical world
well enough to \emph{imagine} its futures and \emph{act} to shape them. The path
runs through a natively multimodal foundation model---one that perceives,
reasons, and generates across all modalities (image, video, text, speech,
action)---organized around a three-stage cognitive loop: Perceive $\rightarrow$
Imagine $\rightarrow$ Act.

Biological intelligence follows that loop without thinking about it: observe the
world, mentally simulate what should happen, then intervene. Current AI breaks
the loop into disconnected pipelines---a perception model here, a planner there,
a controller bolted on the end, each trained by a different team on a different
dataset. I propose building it as one model.

Every major AI capability challenge today---controllable video, computer use,
robotics, world simulation, autonomous driving, creative tools---is a special
case of \emph{multimodal conditional generation}. One architecture subsumes them
all \cite{ref19}. Improvements to the shared backbone benefit every task at once.

The rest of this paper makes that concrete: the cognitive loop, the
mixture-of-transformers architecture that runs it, the staged pretraining recipe,
the self-play flywheel, and the unification of today's siloed problems.

\section{Cognitive Loop: Perceive, Imagine, Act}
Three stages, run in a cycle. Each is inference in the \emph{same} unified
multimodal space---not three separate modules wired together.

\begin{enumerate}
\item \textbf{Perceive.} The model ingests the current state across all
modalities---video, audio, proprioception, language---into a unified
representation. Perception is inference, not a preprocessing step.
\item \textbf{Imagine.} Before acting, the model thinks in multimodal space. Not
just text chain-of-thought, but a \emph{mental video} of the planned trajectory:
predicted proprioception, anticipated contact forces. This multimodal
imagination \emph{is} the world model \cite{ref11,ref12}.
\item \textbf{Act.} The model generates whatever output the task
requires---motor commands, keyboard/mouse, speech---conditioned on the imagined
plan, not on reactive perception alone. The action changes the world, new
percepts arrive, and the loop repeats.
\end{enumerate}

Consider the canonical example: a model deciding to pick up a glass. In the
\emph{Perceive} stage, video, audio, proprioception, and language stream in and
are inferred into one unified representation; the model sees the glass, the
table, and its own hand position, all in the same space. In the \emph{Imagine}
stage, the model rolls out a mental video of the plan: the hand approaching,
fingers closing, the predicted weight shift and contact forces. In the
\emph{Act} stage, conditioned on the imagined plan, the model emits motor
commands that close the grasp; the world changes, fresh percepts arrive, and the
loop returns to Perceive.

The key idea: \textbf{Imagine and Act use the same generative machinery as
Perceive}. Imagining ``the hand approaching, fingers closing, the weight
shifting'' is video + proprioception generation. Acting is action generation.
They are not bolted-on heads---they are the same model generating in different
output modalities. That is what the next section's architecture buys us.

\section{One Architecture: Mixture of Transformers}
If one model has to perceive \emph{and} generate images, video, audio, and
actions, a single dense network creates a problem: modalities interfere. The
statistics of a video frame and a motor command have nothing in common, and
forcing the same feed-forward weights to model both means each capability dilutes
the other.

Start with how the model even sees a video frame or a motor command:
\textbf{everything becomes tokens}. A per-modality encoder maps each input into
one shared token space---text by byte-pair encoding, images and video into VAE
latent patches, audio into neural-codec tokens, actions into discretized or
continuous control tokens. The transformer then sees a single interleaved stream,
each token tagged with its modality.

The architecture I use is a \textbf{mixture of transformers (MoT)}
\cite{ref1}. The trick is to split the network by \emph{modality}:

\begin{itemize}
\item \textbf{Every weight is modality-specific.} Each modality gets its own copy
of the entire transformer block---the query/key/value and output projections, the
feed-forward network, and the RMSNorm layers. Those weights \emph{are} the
experts. A token's modality deterministically selects which expert processes it;
there is no learned router.
\item \textbf{Only the attention \emph{operation} is global.} Each token's Q, K,
and V are produced by its own modality's projections, but attention then runs
over the combined sequence---every token reads all prior context regardless of
modality, so a later video token attends to earlier text and action tokens. The
pattern tracks the objective: autoregressive tokens attend \emph{causally} to the
past, while a block of tokens being generated by diffusion attends
\emph{bidirectionally} within itself, conditioned on that causal context
\cite{ref2,ref3}. The operation has no weights of its own; it is the wire that
makes this \emph{one} model, while every learnable parameter stays
modality-specific.
\end{itemize}

The experts share no weights---only a single representational space they read and
write through global attention; that shared space is the \emph{unified multimodal
space} the loop runs in. Concretely, for token $i$ with modality $m_i$, every
projection carries the modality subscript while attention mixes across the whole
sequence:
\begin{equation}
\begin{aligned}
q_i &= W^{Q}_{m_i}x_i, \quad k_i = W^{K}_{m_i}x_i, \quad v_i = W^{V}_{m_i}x_i \\[4pt]
\alpha_{ij} &= \operatorname{softmax}_j\!\big(q_i^{\top}k_j/\sqrt{d}\big) \\[4pt]
y_i &= \mathrm{FFN}_{m_i}\!\Big(\textstyle\sum_{j}\alpha_{ij}\,v_j\Big)
\end{aligned}
\end{equation}

Every learnable matrix carries the subscript $m_i$. The attention coefficients
$\alpha_{ij}$ are activations, not parameters---computed at runtime---and the
mixing they perform runs globally across modalities, which is exactly what lets
one modality's tokens read another's.

The payoff shows up in the next section: because every weight a token touches
lives in its modality's expert, you can \textbf{freeze an entire expert}---its
attention projections included---while training another, and global attention
still lets the new expert read everything the frozen ones know.

Notice that the attention \emph{operation} never changes---only which modality
experts (and therefore which Q/K/V, output, FFN, and norm weights) are present
and active. Drop video and audio and you have a vision-language understanding
model. Add the action expert and the same backbone becomes an agent.
\textbf{One model, reconfigured by which experts are switched on.}

\section{The Pretraining Recipe: A Staged Curriculum}
The recipe mirrors how we train LLMs, but extends it into multimodal space
\cite{ref19}. Understanding is autoregressive next-token prediction; generation
uses the objective each modality calls for---diffusion for continuous signals
like image, video, audio, and continuous motor control \cite{ref2}---so the same
backbone is autoregressive where it perceives and diffusion where it generates.
Summed over the understanding token positions $\mathcal{U}$ and the generation
positions $\mathcal{G}$, one pretraining loss carries both objectives:
\begin{equation}
\mathcal{L} =
\underbrace{-\sum_{i\in\mathcal{U}}\log p_\theta\big(x_i \mid x_{<i}\big)}_{\text{autoregressive --- understand}}
+ \lambda\,
\underbrace{\sum_{i\in\mathcal{G}}\mathbb{E}_{t,\epsilon}\big\|\,\epsilon - \epsilon_\theta\big(x_i^{(t)},\,t \mid c_i\big)\big\|^2}_{\text{diffusion --- generate}}
\end{equation}
where $c_i$ is the multimodal context attended to by token $i$, and $\lambda$
balances the two objectives, which live on different scales. You do not learn to
perceive, imagine, and act all at once: generation is harder than understanding,
and each generative capability builds on the one before it. So pretraining climbs
a ladder:
\begin{equation*}
\text{MM understanding} \rightarrow \text{image gen.} \rightarrow \text{video / audio gen.} \rightarrow \text{action gen.}
\end{equation*}

Two rules govern the climb:
\begin{itemize}
\item \textbf{Add one expert per stage.} Each new stage introduces a new MoT
generation expert (Section III) on top of everything learned so far.
\item \textbf{Freeze everything earlier.} When training a later generation
component, the earlier components are frozen. Only the newest expert learns.
\end{itemize}

Formally, at stage $s$ only the new expert's parameters $\theta_s$ are trained,
with every earlier expert held fixed:
\begin{equation}
\theta_s^{\star} = \arg\min_{\theta_s}\ \mathcal{L}_s\big(\theta_s \,;\, \underbrace{\theta_1^{\star},\dots,\theta_{s-1}^{\star}}_{\text{frozen}}\big)
\end{equation}

\subsection{Why this order?}
The sequence is not arbitrary---each arrow is a dependency. \emph{Understanding
before generation}: you cannot imagine a coherent future until you can perceive
the present; stage 1 builds the perceptual prior that everything conditions on.
\emph{Image before video}: a video is images over time, so the video expert
reuses the frozen image expert's spatial priors and only has to learn temporal
dynamics---a much smaller lift than learning pixels and motion together.
\emph{Generation before action}: action is conditioned on \emph{imagination}
(Section II), so the model must be able to roll out a predicted
video/proprioception trajectory before it can choose a motor command that
achieves it; the action expert comes last, reading from a frozen world model
\cite{ref23}.

\subsection{The data pyramid}
Staging is also the only practical option, because \textbf{data gets scarcer at
every stage}. The corpus of text, images, video, and audio for
\emph{understanding} is internet-scale. Paired data for high-quality image
\emph{generation} is smaller; video and audio generation smaller still; and
action data---robot trajectories, labeled screen recordings, driving logs---is
the scarcest of all. Each stage trains on roughly an order of magnitude less data
than the one before. So we spend the abundant data first to build the strongest
possible backbone, then climb into the data-poor regimes on top of it.

\subsection{Why freeze?}
Freezing earlier components does three things at once. \emph{No catastrophic
forgetting}: hard-won perception and image priors cannot be eroded by the noisier
gradients of a new, harder generative task; this matters most because of the data
pyramid, since if everything trained jointly the tiny action dataset would be
swamped by---and would perturb---the enormous understanding corpus. \emph{Clean
credit assignment}: the newest expert is the only thing learning, so every
gradient is attributable to it, and nothing else can drift to ``explain away''
its errors. \emph{Cheaper training}: a frozen expert is just a forward pass---no
optimizer state, no backward pass through its weights, so each later stage is far
cheaper than training the whole stack jointly. Because every weight---attention
projections included---lives inside a modality expert, freezing an expert freezes
its Q/K/V too. The global attention operation has no parameters of its own; it
simply lets each newly added expert read and write through the frozen experts'
representations.

\subsection{The learning-rate schedule}
All of this rides on a single learning-rate schedule with three phases: a short
\textbf{warm-up} that ramps the LR to its peak, a long \textbf{constant-LR} (CLR)
plateau that does the bulk of the work, and a \textbf{ramp-down} that anneals the
LR toward zero. The staged curriculum maps directly onto it. The first \textbf{50\%
of the CLR plateau is understanding only}---the data-rich foundation---and the
generation stages (image, then video/audio, then action) are layered into the
second half. The ramp-down is reserved for \textbf{high-quality data across every
modality}, polishing the whole stack as the LR decays to zero.

The staged-and-frozen recipe turns one terrifying joint-training problem into
four tractable ones. Each stage asks a single, well-posed question---\emph{given
everything I already know, how do I generate this one new modality?}---and answers
it without disturbing the rest.

\section{Post-training and the Physics Flywheel}
Pretraining learns the prior over how the multimodal world works. Two stages turn
that prior into a capable agent:
\begin{itemize}
\item \textbf{Post-training (SFT).} Supervised fine-tuning on curated Perceive
$\rightarrow$ Imagine $\rightarrow$ Act demonstrations, with explicit multimodal
reasoning traces---the imagination made legible.
\item \textbf{RL via self-play in simulation} with automatic reward signals. This
is where the model becomes \emph{self-improving}.
\end{itemize}

In simulation, the model runs its full loop and grades itself automatically on
three questions. \emph{Did what I imagined actually happen?}---imagination
accuracy. \emph{Did I succeed?}---task completion. \emph{Does my prediction
violate physics?}---consistency, checked against engine constraints. Call it
\emph{RL from physics feedback} (RLPF).

Self-play maximizes expected return under a composite reward that folds these
together:
\begin{equation}
\max_{\theta}\ \mathbb{E}_{\tau\sim\pi_\theta}\!\Big[\textstyle\sum_t R_t\Big],
\qquad R = \alpha\,R_{\text{task}} + \beta\,R_{\text{imag}} + \gamma\,R_{\text{phys}}
\end{equation}
with $R_{\text{imag}}$ the imagination accuracy and $R_{\text{phys}}$ the
physics-consistency term (RLPF).

Trace one lap to see why it \emph{compounds}. A better world model makes the
model's imagination more accurate, so it can plan against a mental simulation it
trusts---better plans. Better plans yield better actions, which succeed more
often. Those successful rollouts (``this plan led to this outcome'') are exactly
the labeled examples you want, so they become richer training data---which trains
an even better world model. Each lap raises the floor for the next.

\subsection{Why physics, not just text rewards?}
Text models can self-improve this way too---but only where an automatic grader
exists. Math has checkable answers, code has unit tests, and RL from verifiable
rewards has driven much of the recent progress in exactly those domains
\cite{ref29}. Two things limit it. First, \textbf{breadth}: verifiable graders
cover a narrow slice of what we care about, and everything else falls back on
learned proxies---reward models, LLM judges---which are themselves models that can
be reward-hacked and carry no ground truth about the physical world. Second,
\textbf{what the signal improves}: a unit test sharpens a policy, but it does not
teach the model how the world works. Physics fixes both. It is a broad, grounded
grader across the entire embodied space---did the glass actually lift? did the
motion obey gravity?---and that signal flows back into the world model itself.
Physics is to physical reality what unit tests are to code.

The catch is that the whole loop runs \emph{inside} a simulation, so its ceiling
is \textbf{simulation fidelity}. Train against a simulator that gets physics wrong
and the model learns from a lie---the rewards are noise and nothing compounds. And
neither obvious option is enough on its own: a hand-built physics engine is
accurate but cannot render the visual richness of the real world, while real video
is realistic but you cannot act inside it. The essential ingredient is therefore a
\emph{learned} world simulator---trained on real video for realism and
interactivity \cite{ref13}, and corrected by physics engines for physical
consistency \cite{ref20}. Low fidelity and the flywheel slips; high fidelity and
each cycle multiplies the last. The whole strategy lives or dies on the quality of
the imagined world.

\section{One Model, Many Problems}
Here is the payoff of building it as one model. Every major AI capability
challenge is a different \emph{route} through the same architecture---a particular
choice of input modalities and output modality. Familiar combinations name an
established task; the rest are potential applications---routes the same model
could serve, with no household name yet.

Each is the same conditional distribution---choose which modality streams are
given and which are generated:
\begin{equation}
y \sim p_\theta\big(y \mid x_1, x_2, \dots, x_k\big)
\end{equation}

Table~\ref{tab:routes} reads off several familiar tasks as routes through this
one conditional model. Today each of these is a separate model, dataset, and
research community. The Perceive $\rightarrow$ Imagine $\rightarrow$ Act framework
unifies them as different routes through one architecture---and improvements to
the shared backbone benefit every task simultaneously. That is the whole bet:
\textbf{build the multimodal mind once, and every capability comes along for the
ride.}

\begin{table}[t]
\centering
\caption{Today's siloed problems as routes through one multimodal conditional
generator: a choice of input modalities and output modality.}
\label{tab:routes}
\small
\begin{tabular}{@{}lll@{}}
\toprule
inputs (perceive) & output (act) & task \\
\midrule
text & image & text-to-image \\
text & video & text-to-video \\
text, image & video & image animation \\
image, action & video & world simulation \\
video, language & action & robot control (VLA) \\
screen, instruction & action & computer use \\
video, sensors & action & autonomous driving \\
audio, text & speech & speech assistant \\
\bottomrule
\end{tabular}
\end{table}

\section{Related Work}
This essay synthesizes several lines of work. The architecture draws on sparse,
modality-decoupled transformers and unified autoregressive-plus-diffusion models~\cite{ref1,ref2,ref3,ref4,ref5,ref6,ref7,ref8,ref9,ref10};
the cognitive loop and physics flywheel build on the world-models tradition~\cite{ref11,ref12,ref13,ref14,ref15,ref16,ref17,ref18}; and
the physical-AI framing tracks a recent wave of world foundation models~\cite{ref19,ref20,ref21,ref22,ref23,ref24,ref25,ref26,ref27,ref28}. Some of
it---most directly NVIDIA's Cosmos 3 \cite{ref19}---arrived at strikingly similar
conclusions: a two-tower mixture-of-transformers that unifies understanding, world
generation, and action for physical AI. The self-play training loop draws on reinforcement learning from verifiable rewards and modern post-training~\cite{ref29,ref30,ref31,ref32,ref33,ref34}, and on the self-play-and-planning lineage~\cite{ref35,ref36}.

\begin{thebibliography}{99}
\footnotesize
\bibitem{ref1} W. Liang et al., ``Mixture-of-Transformers: A Sparse and Scalable
  Architecture for Multi-Modal Foundation Models,'' TMLR, 2025. arXiv:2411.04996.
\bibitem{ref2} C. Zhou et al., ``Transfusion: Predict the Next Token and Diffuse
  Images with One Multi-Modal Model,'' 2024. arXiv:2408.11039.
\bibitem{ref3} Chameleon Team (Meta FAIR), ``Chameleon: Mixed-Modal Early-Fusion
  Foundation Models,'' 2024. arXiv:2405.09818.
\bibitem{ref4} J. Xie et al., ``Show-o: One Single Transformer to Unify
  Multimodal Understanding and Generation,'' ICLR, 2025. arXiv:2408.12528.
\bibitem{ref5} X. Wang et al. (BAAI), ``Emu3: Next-Token Prediction is All You
  Need,'' 2024. arXiv:2409.18869.
\bibitem{ref6} C. Wu et al. (DeepSeek-AI), ``Janus: Decoupling Visual Encoding for
  Unified Multimodal Understanding and Generation,'' 2024. arXiv:2410.13848.
\bibitem{ref7} X. Chen et al. (DeepSeek-AI), ``Janus-Pro: Unified Multimodal
  Understanding and Generation with Data and Model Scaling,'' 2025. arXiv:2501.17811.
\bibitem{ref8} J. Lu et al. (AI2), ``Unified-IO 2: Scaling Autoregressive
  Multimodal Models with Vision, Language, Audio, and Action,'' CVPR, 2024.
  arXiv:2312.17172.
\bibitem{ref9} T. Li et al., ``Autoregressive Image Generation without Vector
  Quantization (MAR),'' NeurIPS, 2024. arXiv:2406.11838.
\bibitem{ref10} C. Deng et al. (ByteDance Seed), ``Emerging Properties in Unified
  Multimodal Pretraining (BAGEL),'' 2025. arXiv:2505.14683.
\bibitem{ref11} D. Ha and J. Schmidhuber, ``World Models,'' NeurIPS, 2018.
  arXiv:1803.10122.
\bibitem{ref12} D. Hafner et al., ``Mastering Diverse Domains through World Models
  (DreamerV3),'' Nature, 2025. arXiv:2301.04104.
\bibitem{ref13} J. Bruce et al., ``Genie: Generative Interactive Environments,''
  ICML, 2024. arXiv:2402.15391.
\bibitem{ref14} V. Micheli et al., ``Transformers are Sample-Efficient World
  Models (IRIS),'' ICLR, 2023. arXiv:2209.00588.
\bibitem{ref15} E. Alonso et al., ``Diffusion for World Modeling: Visual Details
  Matter in Atari (DIAMOND),'' NeurIPS, 2024. arXiv:2405.12399.
\bibitem{ref16} A. Hu et al. (Wayve), ``GAIA-1: A Generative World Model for
  Autonomous Driving,'' 2023. arXiv:2309.17080.
\bibitem{ref17} S. Gao et al., ``Vista: A Generalizable Driving World Model with
  High Fidelity and Versatile Controllability,'' NeurIPS, 2024. arXiv:2405.17398.
\bibitem{ref18} A. Bar et al. (Meta FAIR / NYU), ``Navigation World Models,''
  CVPR, 2025. arXiv:2412.03572.
\bibitem{ref19} NVIDIA, ``Cosmos 3: Omnimodal World Models for Physical AI,''
  2026. arXiv:2606.02800.
\bibitem{ref20} NVIDIA, ``Cosmos World Foundation Model Platform for Physical
  AI,'' 2025. arXiv:2501.03575.
\bibitem{ref21} NVIDIA, ``World Simulation with Video Foundation Models for
  Physical AI (Cosmos-Predict2.5),'' 2025. arXiv:2511.00062.
\bibitem{ref22} NVIDIA, ``Cosmos-Reason1: From Physical Common Sense to Embodied
  Reasoning,'' 2025. arXiv:2503.15558.
\bibitem{ref23} NVIDIA, ``GR00T N1: An Open Foundation Model for Generalist
  Humanoid Robots,'' 2025. arXiv:2503.14734.
\bibitem{ref24} Google Robotics, ``RT-1: Robotics Transformer for Real-World
  Control at Scale,'' 2022. arXiv:2212.06817.
\bibitem{ref25} Google DeepMind, ``RT-2: Vision-Language-Action Models Transfer
  Web Knowledge to Robotic Control,'' CoRL, 2023. arXiv:2307.15818.
\bibitem{ref26} M. J. Kim et al. (Stanford), ``OpenVLA: An Open-Source
  Vision-Language-Action Model,'' CoRL, 2024. arXiv:2406.09246.
\bibitem{ref27} K. Black et al. (Physical Intelligence), ``$\pi_0$: A
  Vision-Language-Action Flow Model for General Robot Control,'' 2024.
  arXiv:2410.24164.
\bibitem{ref28} Octo Model Team, ``Octo: An Open-Source Generalist Robot
  Policy,'' RSS, 2024. arXiv:2405.12213.
\bibitem{ref29} DeepSeek-AI, ``DeepSeek-R1: Incentivizing Reasoning Capability in
  LLMs via Reinforcement Learning,'' Nature, 2025. arXiv:2501.12948.
\bibitem{ref30} L. Ouyang et al. (OpenAI), ``Training Language Models to Follow
  Instructions with Human Feedback (InstructGPT),'' NeurIPS, 2022. arXiv:2203.02155.
\bibitem{ref31} H. Lightman et al. (OpenAI), ``Let's Verify Step by Step,'' ICLR,
  2024. arXiv:2305.20050.
\bibitem{ref32} Z. Shao et al. (DeepSeek-AI), ``DeepSeekMath: Pushing the Limits
  of Mathematical Reasoning (GRPO),'' 2024. arXiv:2402.03300.
\bibitem{ref33} N. Lambert et al. (AI2), ``T\"ulu 3: Pushing Frontiers in Open
  Language Model Post-Training,'' 2024. arXiv:2411.15124.
\bibitem{ref34} Kimi Team (Moonshot AI), ``Kimi k1.5: Scaling Reinforcement
  Learning with LLMs,'' 2025. arXiv:2501.12599.
\bibitem{ref35} D. Silver et al. (DeepMind), ``Mastering Chess and Shogi by
  Self-Play with a General Reinforcement Learning Algorithm (AlphaZero),''
  Science, 2018. arXiv:1712.01815.
\bibitem{ref36} J. Schrittwieser et al. (DeepMind), ``Mastering Atari, Go, Chess
  and Shogi by Planning with a Learned Model (MuZero),'' Nature, 2020.
  arXiv:1911.08265.
\end{thebibliography}

\end{document}
